% --- Kapitola 2: Súčasný stav problematiky ---

\subsection{Definícia a význam odporúčacích systémov}
S prudkým rozvojom internetových technológií a inteligentných mobilných zariadení čelí moderný používateľ fenoménu informačného preťaženia. V prostredí s takmer neobmedzeným množstvom digitálneho obsahu a služieb sa odporúčacie systémy (angl. \textit{Recommender Systems}) stali kľúčovým nástrojom na personalizovanú filtráciu informácií. 

Ich primárnou úlohou je predpovedať preferencie používateľa a 
zredukovať rozsiahly priestor dostupných položiek na úzky výber
relevantných možností, čím sa zvyšuje celková spokojnosť 
a efektivita interakcie s platformou. Súčasný výskum sa čoraz 
viac orientuje na spracovanie implicitných dát. Ide o
behaviorálne vzorce, ako sú kliknutia (\textit{click-stream}), 
história navštívených miest či čas strávený interakciou s 
konkrétnym objektom \cite{electronics11010141}.

\subsubsection{Obsahové filtrovanie (Content-Based Filtering)}
Tento prístup je založený na analýze špecifických atribútov položiek. Systém identifikuje vlastnosti objektov (napr. textové informácie, kategórie alebo tagy), o ktoré používateľ v minulosti prejavil záujem, a na ich základe generuje odporúčania pre podobné položky \cite{electronics11010141}.
\begin{itemize}
    \item \textbf{Princíp:} Analýza atribútov a metadát položiek (napr. žáner, kategória) s cieľom nájsť podobnosť s preferenciami používateľa.
    \item \textbf{Využitie:} Efektívne najmä pri službách pracujúcich s textovými dátami, hudbou alebo e-commerce \cite{electronics11010141}.
    \item \textbf{Limity:} Nízka miera diverzity (angl. \textit{serendipity}) – používateľovi nie sú ponúkané nové, neobjavené kategórie obsahu, čo vedie k efektu informačnej bubliny.
\end{itemize}

\subsubsection{Kolaboratívne filtrovanie (Collaborative Filtering)}
Kolaboratívne filtrovanie predpovedá vkus používateľa na základe analýzy správania a hodnotení celej komunity používateľov (tzv. ,,wisdom of the crowd``). 

\begin{itemize}
    \item \textbf{Memory-based:} Hľadá podobnosť priamo medzi používateľmi (\textit{User-based}) alebo položkami (\textit{Item-based}) pomocou algoritmov ako KNN alebo korelačných koeficientov.
    \item \textbf{Model-based:} Využíva pokročilé štatistické metódy a techniky strojového učenia, ako sú zhlukovanie (\textit{Clustering}), \textit{Singular Value Decomposition} (SVD) alebo \textit{Principal Component Analysis} (PCA).
    \item \textbf{Hlavné nedostatky:} Zdroj \cite{electronics11010141} definuje tri kritické obmedzenia:
    \begin{itemize}
        \item \textit{Cold Start:} Problém s odporúčaním položiek pre nových používateľov bez histórie interakcií.
        \item \textit{Sparsity (Riedkosť dát):} Situácia, keď je v systéme príliš málo hodnotení na to, aby bolo možné spoľahlivo vypočítať podobnosť.
        \item \textit{Gray Sheep (Sivé ovce):} Používatelia s veľmi unikátnym vkusom, ktorí nevykazujú podobnosť so žiadnou väčšou skupinou používateľov, čo sťažuje generovanie odporúčaní.
    \end{itemize}
\end{itemize}

\subsubsection{Hybridné systémy}
Hybridné modely kombinujú viaceré techniky s cieľom eliminovať slabiny jednotlivých prístupov, najmä problém riedkosti dát (\textit{sparsity}) a studeného štartu. Podľa Burkeho klasifikácie \cite{electronics11010141} existuje sedem základných metód hybridizácie, ktoré sú uvedené v tabuľke \ref{tab:hybrid_types}.
\begin{table}[ht]
\centering
\caption{Typy hybridných odporúčacích systémov podľa Burkeho \cite{electronics11010141}}
\label{tab:hybrid_types}
\begin{tabular}{|l|p{9cm}|}
\hline
\textbf{Metóda} & \textbf{Princíp fungovania} \\ \hline
Weighted & Skóre z rôznych techník sa kombinuje pomocou numerických váh. \\ \hline
Switching & Systém prepína medzi technikami v závislosti od kontextu. \\ \hline
Mixed & Odporúčania z viacerých zdrojov sú zobrazené súčasne. \\ \hline
Feature Combination & Atribúty z jedného zdroja slúžia ako vstupné črty pre druhý model. \\ \hline
Feature Augmentation & Výstup jednej techniky sa stáva parametrom pre ďalšie spracovanie. \\ \hline
Cascade & Postupné zjemňovanie poradia (jeden systém vytvorí rebríček, druhý ho upresní). \\ \hline
Meta-level & Celý model vygenerovaný jednou technikou slúži ako vstup pre druhú. \\ \hline
\end{tabular}
\end{table}
Hlavným cieľom väčšiny štúdií zaoberajúcich sa hybridnými modelmi po roku 2010 je integrácia tzv. vedľajších informácií (\textit{side information}), ako sú sociálne väzby medzi používateľmi alebo kontextové dáta, ktoré dopĺňajú chýbajúce explicitné hodnotenia. Týmto spôsobom hybridné systémy dosahujú vyššiu presnosť a robustnosť v porovnaní so samostatnými modelmi.

\subsection{Techniky spracovania dát v odporúčacích systémoch}
Pre efektívne fungovanie vyššie spomínaných modelov sa využívajú špecifické techniky dolovania dát (Data Mining):

\begin{enumerate}
    \item \textbf{Text Mining:} Analýza neštruktúrovaných textov pomocou sémantickej analýzy a váhovania pojmov (napr. TF-IDF). Slúži na pochopenie obsahu podujatí a preferencií používateľa.
    \item \textbf{K-Nearest Neighbor (KNN):} Algoritmus slúžiaci na klasifikáciu položiek alebo používateľov na základe ich blízkosti v metrickom priestore (geografickom alebo vkusovom).
    \item \textbf{Clustering (Zhlukovanie):} Automatické rozdeľovanie dát do skupín so spoločnými črtami. Využíva sa najmä pri segmentácii používateľov alebo kategorizácii lokalít.
    \item \textbf{Matrix Factorization:} Pokročilá technika rozkladu matíc, ktorá transformuje interakcie používateľov do priestoru latentných faktorov, čím rieši problém riedkosti dát (Sparsity).
    \item \textbf{Neural Networks:} Modely hlbokého učenia schopné modelovať nelineárne vzťahy a spracovávať komplexné dáta zo sociálnych sietí či senzorov mobilných zariadení.
\end{enumerate}

\subsection{Metriky hodnotenia kvality odporúčaní}
Na overenie úspešnosti algoritmov a kvantifikáciu kvality odporúčaní sa využívajú metriky založené na porovnaní predpovedí systému so skutočnými preferenciami používateľa. Základným nástrojom pre tento výpočet je tzv. konfúzna matica (\textit{Confusion Matrix}), ktorá klasifikuje výsledky do štyroch kategórií podľa tabuľky \ref{tab:confusion_matrix} \cite{electronics11010141}.

\begin{table}[ht]
\centering
\caption{Konfúzna matica pre odporúčacie systémy \cite{electronics11010141}}
\label{tab:confusion_matrix}
\begin{tabular}{|l|c|c|}
\hline
 & \textbf{Položka relevantná} & \textbf{Položka nerelevantná} \\ \hline
\textbf{Odporúčaná} & True Positive (TP) & False Positive (FP) \\ \hline
\textbf{Neodporúčaná} & False Negative (FN) & True Negative (TN) \\ \hline
\end{tabular}
\end{table}



Na základe týchto hodnôt definujeme nasledujúce kvalitatívne indikátory:

\begin{itemize}
    \item \textbf{Presnosť (Precision):} Pomer relevantných odporúčaných položiek k celkovému počtu odporúčaní. Vyjadruje mieru úspešnosti systému pri výbere položiek.
    $$Precision = \frac{TP}{TP + FP}$$

    \item \textbf{Úplnosť (Recall):} Pomer úspešne odporúčaných položiek k celkovému počtu všetkých relevantných položiek, ktoré mal systém nájsť. Recall je identický s metrikou TPR (\textit{True Positive Rate}) \cite{electronics11010141}.
    $$Recall = \frac{TP}{TP + FN}$$

    \item \textbf{F-Measure:} Harmonický priemer presnosti a úplnosti. Táto metrika sa využíva najmä kvôli faktu, že medzi Precision a Recall existuje vzťah nepriamej úmery (\textit{trade-off}).
    $$F = 2 \cdot \frac{Precision \cdot Recall}{Precision + Recall}$$

    \item \textbf{Presnosť (Accuracy):} Pomer všetkých správnych predpovedí (TP + TN) k celkovému počtu prípadov. Zdroj \cite{electronics11010141} však upozorňuje, že pri nevyvážených dátach (\textit{unbalanced data}) môže byť táto metrika zavádzajúca.

    \item \textbf{RMSE (Root Mean Squared Error):} Index využívaný na hodnotenie presnosti predikcie číselného hodnotenia. Meria smerodajnú odchýlku rozdielov medzi predpovedaným a skutočným hodnotením.

    \item \textbf{ROC a AUC:} ROC krivka vizualizuje vzťah medzi mierou falošných pozitív (FPR) a mierou skutočných pozitív (TPR). Keďže krivka je grafické vyjadrenie, pre kvantitatívne porovnanie sa využíva plocha pod touto krivkou – \textbf{AUC} (\textit{Area Under the Curve}). Hodnota AUC nad 0,8 sa považuje za indikátor vysokej presnosti modelu \cite{electronics11010141}.
\end{itemize}

%https://www.mdpi.com/2079-9292/11/1/141


